\documentclass[a4paper]{article}

%%%% Packages %%%%
\usepackage{amsfonts}
\usepackage{amsmath}
\usepackage{amssymb}
\usepackage{amsthm}
\usepackage{bbm}
\usepackage{booktabs}
\usepackage{calc}
\usepackage{caption}
\usepackage{centernot}
\usepackage{circuitikz}
\usepackage{dsfont}
\usepackage{esint}
\usepackage{esvect}
\usepackage{float}
\usepackage{fullpage}
\usepackage{graphicx}
\usepackage{hhline}
\usepackage{hyperref}
\usepackage{ifthen}
\usepackage{lipsum}
\usepackage{marvosym}
\usepackage{mathrsfs}
\usepackage{mathtools}
\usepackage{microtype}
\usepackage{multicol}
\usepackage{parskip}
\usepackage{siunitx}
\usepackage{svg}
\usepackage{tikz}
\usepackage{tikz-cd}
\usepackage{titlesec}
\usepackage{titling}
\usepackage{wrapfig}
\usepackage{xcolor}
\usepackage[inline]{enumitem}
\usepackage[makeroom]{cancel}
\usepackage[most]{tcolorbox}
\usetikzlibrary{angles}
\usetikzlibrary{arrows.meta}
\usetikzlibrary{bending}
\usetikzlibrary{calc}
\usetikzlibrary{patterns}
\usetikzlibrary{decorations.pathmorphing}
\usetikzlibrary{decorations.markings}

\setsvg{inkscapeexe="C:/Program Files/Inkscape/bin/inkscape.exe"}

\hypersetup{
	pdfborder={0 0 0}
}
\makeatletter
\renewcommand{\@makefnmark}{\hbox{\textsuperscript{\color{red}\normalfont\@thefnmark}}}
\makeatother
%%%% Packages %%%%
%%%% Custom Commands %%%%

\DeclareMathOperator{\Div}{div}

% Define a problem counter that resets with the section
\newcounter{problem}[section]

% Redefine the problem numbering to include section and problem numbers
\renewcommand{\theproblem}{\arabic{problem}}

% Define the Problem command - BG and ENG versions
\newcommand{\Problem}{\refstepcounter{problem}\fbox{\textbf{Problem \theproblem:}} }
\newcommand{\ProblemBG}{\refstepcounter{problem}\fbox{\textbf{Задача \theproblem:}} }

% Solution Definition
\newcommand{\Solution}[1]{\textbf{Solution\ifthenelse{\equal{#1}{}}{}{ #1}:}}
\newcommand{\SolutionBG}[1]{\textbf{Решение\ifthenelse{\equal{#1}{}}{}{ #1}:}}

% This forces LaTeX to reserve vertical space as if every line had a fraction - use in \begin{aligned} .. \end{aligned}
\newcommand{\rowstrut}{\vphantom{\dfrac{1}{2}}}
%%%% Custom Commands %%%%
% Resistor
\tikzset{
	resistor/.style={
		draw=none,
		decorate,
		decoration={
			markings,
			mark=at position 0.5 with {
				% Connectors
				\draw[line] (-\pgfdecoratedcompleteddistance,0) -- (-0.75,0);
				\draw[line] (0.75,0) -- (\pgfdecoratedcompleteddistance,0);
				% Resistor
				\draw[line] (-0.75, 0.35) rectangle (0.75, -0.35);
				\node[above, transform shape] at (0, 0.35) {$R$};
			}
		}
	}
}

% Electrical battery (POSSIBLY FIX WITH RESISTOR CODE FOR OLYMPIAD 22)
\tikzset{
	battery/.style={
		draw=none,
		decorate,
		decoration={
			markings,
			mark=at position 0.5 with {
				% Connectors
				\draw[line] (-\pgfdecoratedcompleteddistance,0) -- (-4,0);
				\draw[line] (4,0) -- (\pgfdecoratedcompleteddistance,0);
				% Battery
				\draw[line] (-4, -20) -- (-4, 20);
				\draw[line] (4, -10) -- (4, 10);
				% Polarity
				\node[above] at (-4, -15) {$+$};
				\node[below] at (4, -15) {$-$};
			}
		}
	}
}

% Electrical switch
\tikzset{
	switch/.style={
		draw=none,
		decorate,
		decoration={
			markings,
			mark=at position 0.5 with {
				% Connectors
				\draw[line] (-\pgfdecoratedcompleteddistance,0) -- (-15,0);
				\draw[line] (15,0) -- (\pgfdecoratedcompleteddistance,0);
				% Switch
				\draw[line] (-12.5,0) circle (2.5);
				\draw[line] (12.5, 0) circle (2.5);
				\draw[line] (-10,0) -- (10,7);
			}
		}
	}
}

% Springs or inductors
\tikzset{
	spring/.style={
		decorate,
		decoration={
			coil,
			aspect=0.6,
			segment length=7pt,
			amplitude=6pt,
			pre length=5pt,
			post length=5pt
		},
		line width=1.5
	}
}

% A point for electrical terminals (radius = inner sep + line width / 2)
% USAGE: \node[terminal, xshift=-2.5pt] at (1, 1) {};
\tikzset{
	terminal/.style={
		circle,
		draw=black,
		line width=1.5pt,
		inner sep=1.75pt,
		append after command={
			\pgfextra{
				\edef\terminalname{\tikzlastnode}
				\draw[black, line width=1.5pt]
				([shift={(-4.5pt,-4.5pt)}]\terminalname.center) --
				([shift={(-1.94pt,-1.94pt)}]\terminalname.center);
				\draw[black, line width=1.5pt]
				([shift={(1.94pt,1.94pt)}]\terminalname.center) --
				([shift={(4.5pt,4.5pt)}]\terminalname.center);
			}
		}
	}
}

% Just a point for junctions
\tikzset{
	point/.style={
		circle,
		fill=black,
		inner sep=1.5pt
	}
}

% Standardized line
\tikzset{
	line/.style={
		line width=1.5pt,
	}
}


%%% THIS IS A HELPER GRID FOR TIKZ PICTURES %%%

%\begin{center}
%	\begin{tikzpicture}
%		
%		% Grid
%		\draw[help lines, step=1] (-2,-2) grid (12,12);
%		
%		% Axes
%		\draw[->, line width=1] (-2,0) -- (12.5,0) node[right] {$x$};
%		\draw[->, line width=1] (0,-2) -- (0,12.5) node[above] {$y$};
%		
%		% Axis labels
%		\foreach \x in {-2,-1,1,2,...,12}
%		\node[below] at (\x,0) {\small $\x$};
%		
%		\foreach \y in {-2,-1,1,2,...,12}
%		\node[left] at (0,\y) {\small $\y$};
%		
%		% Your drawing
%		\draw[spring] (0,0) -- (10,10);
%		
%	\end{tikzpicture}
%\end{center}



%%% ROTATING ARROW HEADS %%%

% To rotate {Stealth} in place use the flex' command:
%\draw[
%line,
%-{Stealth[length=7pt,flex'=0.9]}] 
%(175,186.5) arc (90:270:12.5);



%%% EXAMPLE HATCHING FOR SURFACES %%%
%\begin{scope}
%	\clip (350,100) rectangle (425,275);
%	
%	\foreach \x in {175,183,...,425} {
%		\draw[line width=0.75pt]
%		(\x,100) -- ++(175,+175);
%	}
%\end{scope}



%%% EXAMPLE ANGLE ARC $$$
% Angles
%\draw[line]
%(7,3)
%arc[start angle=-180, end angle=0, radius=2];
% Pattern Info

\tikzset{
	pattern size/.store in=\mcSize, 
	pattern size = 5pt,
	pattern thickness/.store in=\mcThickness, 
	pattern thickness = 0.3pt,
	pattern radius/.store in=\mcRadius, 
	pattern radius = 1pt}
\makeatletter
\pgfutil@ifundefined{pgf@pattern@name@_bb9kre1tw}{
	\pgfdeclarepatternformonly[\mcThickness,\mcSize]{_bb9kre1tw}
	{\pgfqpoint{0pt}{-\mcThickness}}
	{\pgfpoint{\mcSize}{\mcSize}}
	{\pgfpoint{\mcSize}{\mcSize}}
	{
		\pgfsetcolor{\tikz@pattern@color}
		\pgfsetlinewidth{\mcThickness}
		\pgfpathmoveto{\pgfqpoint{0pt}{\mcSize}}
		\pgfpathlineto{\pgfpoint{\mcSize+\mcThickness}{-\mcThickness}}
		\pgfusepath{stroke}
}}
\makeatother

\usepackage[english]{babel}

\usepackage[
left=2.5cm,
right=2.5cm,
top=3cm,
bottom=2.5cm
]{geometry}

\setlength{\droptitle}{-7em}
\title{\Huge Physics Olympiad L2-26}
\author{BigChunguz24 - Last Edit 16/10/2026}
\date{}

\begin{document}
	\maketitle
	
	\Problem
	\textbf{\underline{Dielectric Ball}}
	% Irodov 3.96, 3.98
	\\
	In this question we study a uniform dielectric ball of radius $R$ and relative permittivity $\varepsilon$ in two settings:
	\begin{description}[leftmargin=0pt, labelwidth=0pt,]
		\item[Setting 1:] The dielectric ball is placed in vacuum in a uniform external electric field of strength $\vv{E}_0$.
		\begin{description}[
			leftmargin=!,
			labelwidth=\widthof{$\bullet$},
			labelindent=\widthof{$\bullet\bullet$}
			]
			\item[$\bullet$] Find the electric field strength $\vv{E}$ inside the ball and the polarization $\vv{P}$ of the dielectric; 
		\end{description}
		\begin{description}[
			leftmargin=!,
			labelwidth=\widthof{$\bullet$},
			labelindent=\widthof{}
			]
			\item[] A ball that is polarized in this way may be represented as a result of a small rigid shift $d \ll R$ of all positive charges of the dielectric relative to all the negative ones. It can be assumed that there are $n$ atoms per unit volume of the dielectric, each with atomic number $Z$.
		\end{description}
		\begin{description}[
			leftmargin=!,
			labelwidth=\widthof{$\bullet$},
			labelindent=\widthof{$\bullet\bullet$}
			]
			\item[$\bullet$] Justify why the electric field outside of the ball behaves like a dipole at the centre of the ball;
			\item[$\bullet$] Find an expression for the total electric potential outside of the ball;
			\item[$\bullet$] Find an expression for the total electric field outside of the ball;
			\item[$\bullet$] Find an expression for the small shift $d$;
		\end{description}
		\item[Setting 2:] Three such dielectric balls are fixed on the vertices of an equilateral triangle of side $a \gg R$. The configuration is then placed in a uniform external electric field of strength $\vv{E}_0$ so that the field is perpendicular to one of the sides of the triangle. Assume that $\vv{E}_0$ acts in the plane of the triangle and that enough time has passed for the system to reach equilibrium.
		\begin{description}[
			leftmargin=!,
			labelwidth=\widthof{$\bullet$},
			labelindent=\widthof{$\bullet\bullet$}
			]
			\item[$\bullet$] Taking the potential energy to be zero when the three balls are infinitely separated from each other, find the potential energy of this configuration;
		\end{description}
	\end{description}
	\emph{\underline{Note:}} You should not use the dipole trick for the first part of ``Setting 1''. 
	
	\Problem 
	\textbf{\underline{Fresnel Mirrors}}
	% Irodov 5.71
	\\
	The below figure illustrates the interference experiment with Fresnel mirrors $M_1$ and $M_2$. The angle between the mirrors is $\varphi = 12'$, the distances from their intersection line to the narrow slit $S$ and the screen are equal to $r = 10\,\mathrm{cm}$ and $b = 130\,\mathrm{cm}$ respectively. The wavelength of light is $\lambda = 0{,}55\,\mu\mathrm{m}$. Find:
	\begin{description}[leftmargin=!, labelwidth=\widthof{$\bullet$}]
		\item[$\bullet$] The width of a fringe on the screen and the number of possible maxima;
		\item[$\bullet$] The shift of the interference pattern on the screen when the slit is displaced by $\delta l = 1\,\mathrm{mm}$ along the arc of radius $r$ with centre at the point of intersection of the mirrors;
		\item[$\bullet$] At what maximum width $\delta_{max}$ of the slit the interference fringes on the screen are still observed sufficiently sharp.
	\end{description}
	
	\begin{center}
	\begin{tikzpicture}
		% Screen
		\draw[line] (10, 1) -- (10, 7);
		
		% Distance to mirror and to source
		\draw[line, dashed] (5, 4) -- (10, 4)
		node[midway, above] {$b$};
		\draw[line, dashed] (5, 4) -- (4, 7)
		node[midway, right] {$r$};
		
		% Source
		\filldraw[fill=black, draw=black] (5,4) circle (2.8pt);
		\node[point] at (4, 7) {};
		\node[left] at (4, 7) {$S$};
		
		% Mirrors
		\draw[line] (5, 4) -- (8, 2);
		\draw[line, dashed] (5, 4) -- (1.7, 6.2);
		\draw[line] (5, 4) -- (2.5, 6.5);
		
		\node[above] at (2.5, 6.5) {$M_1$};
		\node[below] at (8, 2) {$M_2$};
		
		% Phantom sources
		\draw[line, dashed] (5, 4) -- (0.5, 4);
		
		\node[point] at (1, 4.2) {};
		\node[point] at (1, 3.8) {};
		
		\node[above] at (1, 4.2) {$S'$};
		\node[below] at (1, 3.8) {$S''$};
		
		% Hatching
		\begin{scope}
			\clip (10,1) rectangle (10.3,7);
			
			\foreach \y in {1,1.3,...,8} {
				\draw[line width=0.75pt]
				(10,\y) -- ++(1,1);
			}
		\end{scope}
	\end{tikzpicture}
\end{center}
	
	\Problem
	\textbf{\underline{Polarization}}
	% Irodov 5.178
	\\
	What kind of polarization has a plane electromagnetic wave if the projections of the vector $\vv{E}$ on the $x$ and $y$ axes are perpendicular to the propagation direction and are defined as:
	\begin{description}
		\item[$\bullet$] $E_x = E \cos(\omega t - kz)$ and $E_y = E \cos(\omega t - kz)$
		\item[$\bullet$] $E_x = E \cos(\omega t - kz)$ and $E_y = E \cos(\omega t - kz + \pi/4)$
		\item[$\bullet$] $E_x = E \cos(\omega t - kz)$ and $E_y = E \cos(\omega t - kz + \pi/2)$
		\item[$\bullet$] $E_x = E \cos(\omega t - kz)$ and $E_y = E \cos(\omega t - kz + 3\pi/4)$
		\item[$\bullet$] $E_x = E \cos(\omega t - kz)$ and $E_y = E \cos(\omega t - kz + \pi)$
	\end{description}
	
	\Problem
	\textbf{\underline{Efficiency}}
	% Irodov 2.128
	\\
	Using the Clausius inequality, show that all cycles having the same maximum temperature $T_{max}$ and minimum temperature $T_{min}$ are less efficient compared to the Carnot cycle with the same $T_{max}$ and $T_{min}$.
	
	\Problem
	\textbf{\underline{The Reduced Compton Wavelength and Compton Scattering}}
	% Irodov 5.280, 5.281, 5.282 + Original
	\begin{description}[leftmargin=!, labelwidth=\widthof{$\bullet$}]
		\item[$\bullet$] \textbf{The Scale and The Reduced Compton Wavelength}
		\\
		In this problem you could assume that the effective mass of a photon is given by the following relation:
		\begin{equation*}
			\boxed{\hbar \omega = m_{eff} c^2}
		\end{equation*}
		
		\begin{itemize}[leftmargin=0.5cm]
			\item[$\rightarrow$] Compose an expression for a quantity whose dimension is length, using the velocity of light $c$, mass of a particle $m$, and Planck's constant $\hbar$. This expression is the reduced Compton wavelength $\bar{\lambda}$;
			\item[$\rightarrow$] Evaluate $\bar{\lambda}_e$ for an electron and find the frequency $\nu$ of a photon of the same reduced wavelength $\bar{\lambda}_e$;
			\item[$\rightarrow$] Estimate the spatial scale $\Delta x$ to which an electron must be confined for its kinetic energy to become comparable to its rest energy. What qualitatively new process becomes possible below this size?
			\item[$\rightarrow$] In the \textbf{Bohr model of hydrogen} the radius $a_0$ of the ground-state orbit is called the
			\textbf{Bohr radius}. The ratio of the electron's
			reduced Compton wavelength $\bar{\lambda}_e$ to the
			Bohr radius defines the \textbf{fine-structure constant} $\alpha$:
			\begin{equation*}
				\boxed{\alpha \coloneq \frac{\bar{\lambda}_e}{a_0}}
			\end{equation*}
			Derive an expression for $\alpha$ in terms of $e$, $\varepsilon_0$,
			$\hbar$ and $c$, and evaluate its numerical value;
		\end{itemize}
		\item[$\bullet$] \textbf{Absorption of a Photon}
		\begin{itemize}[leftmargin=0.5cm]
			\item[$\rightarrow$] Prove that a free electron cannot absorb a photon completely;
			\item[$\rightarrow$] A hydrogen atom, initially at rest in its ground state, absorbs a photon and goes to the excited state $n=2$. Find the exact energy $E_{\gamma}$ of the photon and the kinetic
			energy $E_k$ of the recoiling atom. Hence, evaluate the ratio $\eta$ which shows by how much $E_{\gamma}$ exceeds $\Delta E$:
			\begin{equation*}
				\eta = \frac{\left| E_{\gamma} - \Delta E \right| }{\Delta E}
			\end{equation*}
			where $\Delta E$ is the excitation energy of the atom in its rest frame. Any approximations you make should be thoroughly reasoned for and explicitly shown to be appropriate;
		\end{itemize}
		\item[$\bullet$] \textbf{Compton Scattering}
		\begin{itemize}[leftmargin=0.5cm]
			\item[$\rightarrow$] Derive the Compton shift $\Delta \lambda$ for scattering on a free electron at rest;
			\item[$\rightarrow$] Express the energy $E_\gamma'$ of the
			scattered photon in terms of the incident photon energy
			$E_\gamma$ and the scattering angle $\varphi$. Hence, determine
			the kinetic energy $E_k$ of the recoiling electron;
			\item[$\rightarrow$] Determine the maximum possible wavelength
			shift and the corresponding scattering angle. Express the result
			in terms of the reduced Compton wavelength $\bar{\lambda}_e$
			and evaluate it numerically;
			\item[$\rightarrow$] Determine the maximum kinetic energy of
			the recoiling electron. Examine the limiting cases
			$E_\gamma \ll m_e c^2$ and $E_\gamma \gg m_e c^2$ and explain
			their physical significance;
		\end{itemize}
	\end{description}
	
	\Problem
	\textbf{\underline{Resistor Circuit}}
	% Irodov 3.190
	\\
	A DC source of electromotive force $\varepsilon$ and internal
	resistance $R_0$ is connected to three identical resistors, each
	of resistance $R$, as shown in the figure. The two long connecting branches drawn above and below the three resistors each have resistance $r$, while all other connecting wires have negligible resistance.
	\begin{description}
		\item[$\bullet$] Obtain an expression for the total thermal power $P$ dissipated in the three resistors;
		\item[$\bullet$] Determine the value of $R$ in terms of $R_0$ for which $P$ is maximized when $r = R_0$ and find $P_{max}$;
	\end{description}
	
	\begin{center}
	\begin{tikzpicture}
		% Resistors				
		\draw[resistor] (1, 1) -- (4, 1);
		\draw[resistor] (4, 1) -- (7, 1);
		\draw[resistor] (7, 1) -- (10, 1);
		
		% Terminals
		\draw[line] (1, 1) -- (0.5, 1);
		\node[terminal, xshift=-2.5pt] at (0.5, 1) {};
		
		\draw[line] (10, 1) -- (10.5, 1);
		\node[terminal, xshift=2.5pt] at (10.5, 1) {};
		
		% Junctions
		\node[point] at (1, 1) {};
		\node[point] at (4, 1) {};
		\node[point] at (7, 1) {};
		\node[point] at (10, 1) {};
		
		% Connectors
		\draw[line] (1, 1) -- (1, -0.5) -- (7, -0.5) -- (7, 1);
		\draw[line] (4, 1) -- (4, 2.5) -- (10, 2.5) -- (10, 1);
		
		\node[below] at (4, -0.5) {$r$};
		\node[above] at (7, 2.5) {$r$};
	\end{tikzpicture}
\end{center}
	
	\newpage
\vspace*{-5em}
\enlargethispage{\baselineskip}
\newcommand{\tg}{\operatorname{tg}}
\newcommand{\arctg}{\operatorname{arctg}}
\section*{\centering Appendix}
\section*{\centering Useful Mathematical Formulae and Physical Constants}
\emph{\underline{Important}}: Infinite series that are marked with $(*)$ are only valid under the condition that $|x| < 1$.

\begin{minipage}[t]{0.5\textwidth}
	\vspace{0pt}
	\begin{tcolorbox}[title=\textbf{Taylor Expansions}]
		\vspace{-10pt}
		\begin{align*}
			& \sin(x) = x - \frac{x^3}{6} + \frac{x^5}{120} + \mathcal{O}(x^7)
			&&
			\\[1pt]
			& \cos(x) = 1 - \frac{x^2}{2} + \frac{x^4}{24} - \frac{x^6}{720} + \mathcal{O}(x^8)
			&&
			\\[1pt]
			& \tg(x) = x + \frac{x^3}{3} + \frac{2x^5}{15} + \mathcal{O}(x^7)
			&&
			\\[1pt]
			& \arcsin(x) = x + \frac{x^3}{6} + \frac{3x^5}{40} + \mathcal{O}(x^7)
			&&
			\\[1pt]
			& \arctg(x) = x - \frac{x^3}{3} + \frac{x^5}{5} + \mathcal{O}(x^7)
			&&
			\\[1pt]
			& \exp(x) = 1 + x + \frac{x^2}{2} + \mathcal{O}(x^3)
			&&
			\\[1pt]
			& \ln(1+x) = x - \frac{1}{2}x^2 + \mathcal{O}(x^3)
			&& (*)
			\\[1pt]
			& (1-x)^{-1} = 1 + x + x^2 + \mathcal{O}(x^3)
			&& (*)
			\\[1pt]
			& (1+x)^{-1} = 1 - x + x^2 + \mathcal{O}(x^3)
			&& (*)
			\\[1pt]
			& (1+x)^{+\frac{1}{2}} = 1 + \frac{1}{2}x - \frac{1}{8}x^2 + \mathcal{O}(x^3)
			&& (*)
			\\[1pt]
			& (1+x)^{-\frac{1}{2}} = 1 - \frac{1}{2} x + \frac{3}{8} x^2 + \mathcal{O}(x^3)
			&& (*)
			\\[1pt]
			& (1+x)^{+\frac{3}{2}} = 1 + \frac{3}{2}x + \frac{3}{8}x^2 + \mathcal{O}(x^3)
			&& (*)
			\\[1pt]
			& (1+x)^{-\frac{3}{2}} = 1 - \frac{3}{2}x + \frac{15}{8}x^2 + \mathcal{O}(x^3)
			&& (*)
		\end{align*}
		
	\end{tcolorbox}
\end{minipage}
\begin{minipage}[t]{0.5\textwidth}
	\vspace{0pt}
	\begin{tcolorbox}[title=\textbf{Inverse Hyperbolic Functions}]
		\vspace{-10pt}
		\begin{align*}
			& \text{arsinh}(x) = \ln(x + \sqrt{x^2 + 1})
			\\[3.0375pt]
			& \text{arcosh}(x) = \ln(x + \sqrt{x^2 - 1})
			\\[3.0375pt]
			& \text{artgh}(x) = \frac{1}{2} \ln
			\left(
			\frac{1+x}{1-x}
			\right)
			\\[3.0375pt]
			& \text{arctgh}(x) = \frac{1}{2} \ln
			\left(
			\frac{x+1}{x-1}
			\right)
		\end{align*}
		
	\end{tcolorbox}
	\begin{tcolorbox}[title=\textbf{Indefinite Integrals}]
		\vspace{-10pt}
		\begin{align*}
			& \int (1-x^2)^{-\frac{3}{2}} dx = \frac{x}{\sqrt{1-x^2}} + C
			\\[3.0375pt]
			& \int (1+x^2)^{-\frac{3}{2}} dx = \frac{x}{\sqrt{1 + x^2}} + C
			\\[3.0375pt]
			& \int (1-x^2)^{-\frac{1}{2}} dx = \arcsin(x) + C
			\\[3.0375pt]
			& \int (1+x^2)^{-\frac{1}{2}} dx = \text{arsinh}(x) + C
			\\[3.0375pt]
			& \int (1-x^2)^{+\frac{1}{2}} dx = \frac{\arcsin(x) + x \sqrt{1-x^2}}{2} + C
			\\[3.0375pt]
			& \int (1+x^2)^{+\frac{1}{2}} dx = \frac{\text{arsinh}(x) + x \sqrt{1 + x^2}}{2} + C
		\end{align*}
		
	\end{tcolorbox}
\end{minipage}

\begin{minipage}{\textwidth}
	\begin{tcolorbox}[title=\textbf{Useful Physical Constants}]
		\vspace{-5pt}
		\begin{center}
			\begin{tabular}{|l|c|c|c|}
				\hline
				\textbf{Constant} & \textbf{Symbol} & \textbf{Value} & \textbf{Measuring Units} 
				\\
				\hline
				Light speed in vacuum & $c$ & $3{,}000 \times 10^8$ & $\mathrm{m/s}$
				\rule{0pt}{10pt}
				\\[1pt]
				\hline
				Rest mass of electron & $m_e$ & $9{,}109 \times 10^{-31}$ & $\mathrm{kg}$
				\rule{0pt}{10pt}
				\\[1pt]
				\hline
				Rest mass of proton & $m_p$ & $1{,}673 \times 10^{-27}$ & $\mathrm{kg}$
				\rule{0pt}{10pt}
				\\[1pt]
				\hline
				Rest mass of neutron & $m_n$ & $1{,}675 \times 10^{-27}$ & $\mathrm{kg}$
				\rule{0pt}{10pt}
				\\[1pt]
				\hline
				Rest mass of muon & $m_{\mu}$ & $1{,}884 \times 10^{-28}$ & $\mathrm{kg}$
				\rule{0pt}{10pt}
				\\[1pt]
				\hline
				Rest mass of tau-lepton & $m_{\tau}$ & $3{,}168 \times 10^{-27}$ & $\mathrm{kg}$
				\rule{0pt}{10pt}
				\\[1pt]
				\hline
				Elementary charge & $e$ & $1{,}602 \times 10^{-19}$ & $\mathrm{C}$
				\rule{0pt}{10pt}
				\\[1pt]
				\hline
				Planck's constant & $h$ & $6{,}626 \times 10^{-34}$ & $\mathrm{Js}$
				\rule{0pt}{10pt}
				\\[1pt]
				\hline
				Boltzmann's constant & $k$ & $1{,}381 \times 10^{-23}$ & $\mathrm{J/K}$
				\rule{0pt}{10pt}
				\\[1pt]
				\hline
				Gravitational constant & $\gamma$ & $6{,}673 \times 10^{-11}$ & $\mathrm{Nm^2/kg^2}$ 
				\rule{0pt}{10pt}
				\\[1pt]
				\hline
				Dielectric permittivity of vacuum & $\varepsilon_0$ & $8{,}854 \times 10^{-12}$ & $\mathrm{F/m}$
				\rule{0pt}{10pt}
				\\[1pt]
				\hline
				Avogadro's number & $N_A$ & $6{,}022 \times 10^{23}$ & $\mathrm{mol^{-1}}$
				\rule{0pt}{10pt}
				\\[1pt]
				\hline
				Rydberg's constant & $R_{\infty}$ & $1{,}097 \times 10^7$ & $\mathrm{m^{-1}}$
				\rule{0pt}{10pt}
				\\[1pt]
				\hline
				Magnetic permeability of vacuum & $\mu_0$ & $4\pi \times 10^{-7}$ & $\mathrm{H/m}$
				\rule{0pt}{10pt}
				\\[1pt]
				\hline
				Universal gas constant & $R$ & $8{,}314$ & $\mathrm{J / (mol \cdot K)}$
				\rule{0pt}{10pt}
				\\[1pt]
				\hline
			\end{tabular}
		\end{center}
	\end{tcolorbox}
\end{minipage}
	
	% Irodov 3.96, 3.98
	% Irodov 5.71
	% Irodov 5.178
	% Irodov 2.128
	% Irodov 5.280, 5.281, 5.282 + Original
	% Irodov 3.190
\end{document}